\documentclass[10pt,a4paper]{amsart}
\usepackage{amsmath,amssymb,amsthm,graphicx}
\usepackage{verbatim}
\newenvironment{reference}{\comment}{\endcomment}
%\newenvironment{reference}{}{}
\newenvironment{slogan}{\comment}{\endcomment}
\newenvironment{history}{\comment}{\endcomment}
% For commutative diagrams we use Xy-pic
\usepackage[all]{xy}

% We use 2cell for 2-commutative diagrams.
\xyoption{2cell}
% Theorem environments.
%
\theoremstyle{plain}
\newtheorem{theorem}[subsection]{Theorem}
\newtheorem{proposition}[subsection]{Proposition}
\newtheorem{lemma}[subsection]{Lemma}

\theoremstyle{definition}
\newtheorem{definition}[subsection]{Definition}
\newtheorem{example}[subsection]{Example}
\newtheorem{exercise}[subsection]{Exercise}
\newtheorem{situation}[subsection]{Situation}

\theoremstyle{remark}
\newtheorem{remark}[subsection]{Remark}
\newtheorem{remarks}[subsection]{Remarks}

\numberwithin{equation}{subsection}

% Macros
%
\def\lim{\mathop{\mathrm{lim}}\nolimits}
\def\colim{\mathop{\mathrm{colim}}\nolimits}
\def\Spec{\mathop{\mathrm{Spec}}}
\def\Hom{\mathop{\mathrm{Hom}}\nolimits}
\def\Ext{\mathop{\mathrm{Ext}}\nolimits}
\def\SheafHom{\mathop{\mathcal{H}\!\mathit{om}}\nolimits}
\def\SheafExt{\mathop{\mathcal{E}\!\mathit{xt}}\nolimits}
\def\Sch{\mathit{Sch}}
\def\Mor{\mathop{\mathrm{Mor}}\nolimits}
\def\Ob{\mathop{\mathrm{Ob}}\nolimits}
\def\Sh{\mathop{\mathit{Sh}}\nolimits}
\def\NL{\mathop{N\!L}\nolimits}
\def\CH{\mathop{\mathrm{CH}}\nolimits}
\def\proetale{{pro\text{-}\acute{e}tale}}
\def\etale{{\acute{e}tale}}
\def\QCoh{\mathit{QCoh}}
\def\Ker{\mathop{\mathrm{Ker}}}
\def\Im{\mathop{\mathrm{Im}}}
\def\Coker{\mathop{\mathrm{Coker}}}
\def\Coim{\mathop{\mathrm{Coim}}}

% Boxtimes
%
\DeclareMathSymbol{\boxtimes}{\mathbin}{AMSa}{"02}

%
% Macros for moduli stacks/spaces
%
\def\QCohstack{\mathcal{QC}\!\mathit{oh}}
\def\Cohstack{\mathcal{C}\!\mathit{oh}}
\def\Spacesstack{\mathcal{S}\!\mathit{paces}}
\def\Quotfunctor{\mathrm{Quot}}
\def\Hilbfunctor{\mathrm{Hilb}}
\def\Curvesstack{\mathcal{C}\!\mathit{urves}}
\def\Polarizedstack{\mathcal{P}\!\mathit{olarized}}
\def\Complexesstack{\mathcal{C}\!\mathit{omplexes}}
% \Pic is the operator that assigns to X its picard group, usage \Pic(X)
% \Picardstack_{X/B} denotes the Picard stack of X over B
% \Picardfunctor_{X/B} denotes the Picard functor of X over B
\def\Pic{\mathop{\mathrm{Pic}}\nolimits}
\def\Picardstack{\mathcal{P}\!\mathit{ic}}
\def\Picardfunctor{\mathrm{Pic}}
\def\Deformationcategory{\mathcal{D}\!\mathit{ef}}
\usepackage{cleveref}
\crefname{originallemma}{original lemma}{original lemmas}
\Crefname{originallemma}{Original lemma}{Original lemmas}
\makeatletter
\newlabel{lemma-aut-iso-unramified}{{lemma-aut-iso-unramified}{source}{}{sourcepointer-1}{}}
\newlabel{lemma-aut-iso-unramified@cref}{{[originallemma][0][]lemma-aut-iso-unramified}{[1][0][]source}{}{}{}}
\makeatother
\crefname{originallemma}{original lemma}{original lemmas}
\Crefname{originallemma}{Original lemma}{Original lemmas}
\makeatletter
\newlabel{lemma-stabilizer-preserving}{{lemma-stabilizer-preserving}{source}{}{sourcepointer-2}{}}
\newlabel{lemma-stabilizer-preserving@cref}{{[originallemma][0][]lemma-stabilizer-preserving}{[1][0][]source}{}{}{}}
\makeatother
\crefname{originallemma}{original lemma}{original lemmas}
\Crefname{originallemma}{Original lemma}{Original lemmas}
\makeatletter
\newlabel{lemma-universally-injective}{{lemma-universally-injective}{source}{}{sourcepointer-3}{}}
\newlabel{lemma-universally-injective@cref}{{[originallemma][0][]lemma-universally-injective}{[1][0][]source}{}{}{}}
\makeatother
\crefname{originalremark}{original remark}{original remarks}
\Crefname{originalremark}{Original remark}{Original remarks}
\makeatletter
\newlabel{remark-identify-automorphism-groups}{{remark-identify-automorphism-groups}{source}{}{sourcepointer-4}{}}
\newlabel{remark-identify-automorphism-groups@cref}{{[originalremark][0][]remark-identify-automorphism-groups}{[1][0][]source}{}{}{}}
\makeatother
\crefname{originalreference}{original reference}{original references}
\Crefname{originalreference}{Original reference}{Original references}
\makeatletter
\newlabel{spaces-morphisms-lemma-etale-universally-injective-open}{{spaces-morphisms-lemma-etale-universally-injective-open}{source}{}{sourcepointer-5}{}}
\newlabel{spaces-morphisms-lemma-etale-universally-injective-open@cref}{{[originalreference][0][]spaces-morphisms-lemma-etale-universally-injective-open}{[1][0][]source}{}{}{}}
\makeatother
\title{A stabilizer preserving etale morphism is an open immersion}
\author{The Stacks Project Authors}
\date{stacks-project:stacks-morphisms.tex:local-snapshot-20261012; GNU FDL 1.2 or later}
\begin{document}
\maketitle
\noindent\textit{Editorial note.} All delivered mathematical text is the author's own, copied from the pinned Stacks Project chapter source sources/stacks-algebraic-stacks-morphisms/pinned-original-source.tex; no statement or proof step was rewritten, reordered or removed.\par\medskip
\noindent\textit{Editorial note.} Preamble provenance: the Stacks Project keeps one single shared preamble for every chapter; this card takes it directly from the staged sources/stacks-algebraic-stacks-morphisms/preamble.tex as three verbatim spans (lines 8-12, 14-18 and 154-217), so no derived file was written and nothing was added or redefined. The wrapper class and the base packages are supplied by the assembler.\par\medskip
\noindent\textit{Editorial note.} Bibliography provenance: neither the counted statement nor its proof contains a citation (this lane rescans the whole span before writing the card), so the card carries no bibliography and no entry was reconstructed.\par\medskip
\noindent\textit{Editorial note.} Selection provenance: this card was opened under the file-based control contract (CONTROL.md), with the mechanical pre-card gate run on its statement and proof interval before the card was written.\par\medskip
\noindent\textit{Allowed standard prerequisites (editorial).}\begin{itemize}
\item The statement that a morphism whose diagonal is an isomorphism is universally injective. Invoked at the start of the counted proof (source line 10172). Gives universal injectivity from the hypothesis on the diagonal.
\item The combination of stabilizer preservation with unramifiedness of the automorphism group. Invoked inside the counted proof (source lines 10173-10176). Makes the morphism representable by algebraic spaces.
\item The statement that an etale universally injective morphism of algebraic spaces is an open immersion. Invoked near the end of the counted proof (source lines 10177-10179) as a cross-chapter pointer. Turns the two properties into an open immersion.
\end{itemize}\medskip
\ifdefined\hypertarget\hypertarget{sourcepointer-1}{}\fi
\noindent\textit{Source reference.} Original lemma lemma-aut-iso-unramified: sources/stacks-algebraic-stacks-morphisms/pinned-original-source.tex, source lines 10012--10012. Navigation pointer only: the counted text refers to the Stacks item lemma-aut-iso-unramified, which is an existing labeled result of this same chapter and is not part of this card. The pointer carries the target's own label as its number and its definition line as its span; its content is not reproduced here.\par\medskip
\ifdefined\hypertarget\hypertarget{sourcepointer-2}{}\fi
\noindent\textit{Source reference.} Original lemma lemma-stabilizer-preserving: sources/stacks-algebraic-stacks-morphisms/pinned-original-source.tex, source lines 9980--9980. Navigation pointer only: the counted text refers to the Stacks item lemma-stabilizer-preserving, which is an existing labeled result of this same chapter and is not part of this card. The pointer carries the target's own label as its number and its definition line as its span; its content is not reproduced here.\par\medskip
\ifdefined\hypertarget\hypertarget{sourcepointer-3}{}\fi
\noindent\textit{Source reference.} Original lemma lemma-universally-injective: sources/stacks-algebraic-stacks-morphisms/pinned-original-source.tex, source lines 2907--2907. Navigation pointer only: the counted text refers to the Stacks item lemma-universally-injective, which is an existing labeled result of this same chapter and is not part of this card. The pointer carries the target's own label as its number and its definition line as its span; its content is not reproduced here.\par\medskip
\ifdefined\hypertarget\hypertarget{sourcepointer-4}{}\fi
\noindent\textit{Source reference.} Original remark remark-identify-automorphism-groups: sources/stacks-algebraic-stacks-morphisms/pinned-original-source.tex, source lines 4046--4046. Navigation pointer only: the counted text refers to the Stacks item remark-identify-automorphism-groups, which is an existing labeled result of this same chapter and is not part of this card. The pointer carries the target's own label as its number and its definition line as its span; its content is not reproduced here.\par\medskip
\ifdefined\hypertarget\hypertarget{sourcepointer-5}{}\fi
\noindent\textit{Source reference.} Original reference spaces-morphisms-lemma-etale-universally-injective-open: sources/stacks-algebraic-stacks-morphisms/pinned-original-source.tex, source lines 4656--4656. Navigation pointer only: the counted text refers to the Stacks item spaces-morphisms-lemma-etale-universally-injective-open, which belongs to another chapter of the project and is not part of this card. The pointer carries the target's own label as its number because this chapter's staged snapshot does not contain the other chapter's tag table; its content is not reproduced here.\par\medskip
\begin{lemma}
\label{lemma-etale-iso}
Let $f : \mathcal{X} \to \mathcal{Y}$ be a morphism of algebraic stacks.
Assume $f$ is \'etale, $f$ induces an isomorphism
between automorphism groups at points
(Remark \ref{remark-identify-automorphism-groups}),
and for every algebraically closed field $k$ the functor
$$
f : \Mor(\Spec(k), \mathcal{X}) \longrightarrow \Mor(\Spec(k), \mathcal{Y})
$$
is an equivalence. Then $f$ is an isomorphism.
\end{lemma}

\begin{proof}
By Lemma \ref{lemma-universally-injective} we see that $f$ is
universally injective. Combining
Lemmas \ref{lemma-stabilizer-preserving} and
\ref{lemma-aut-iso-unramified}
we see that $f$ is representable by algebraic spaces.
Hence $f$ is an open immersion by Morphisms of Spaces, Lemma
\ref{spaces-morphisms-lemma-etale-universally-injective-open}.
To finish we remark that the condition in the lemma also guarantees
that $f$ is surjective.
\end{proof}

\end{document}
